\begin{proofbox}
	\textbf{\textcolor{CProof}{思路提示}}

	\textbf{利用 $x_1$ 与 $x_2$ 的比值替换，把二元问题化为只含一个变量 $t$ 的不等式；再用导数分别证明两个基础不等式。}

	\medskip
	\textbf{\textcolor{CProof}{正式推导}}
	\begin{description}[style=nextline, leftmargin=2.8em, labelsep=0.5em, itemsep=0.55em]
		\item[\textbf{步骤一（设比值大小）：}]
		      由对称性，不妨设 $x_1>x_2$。令
		      \[
			      t=\frac{x_1}{x_2},
		      \]
		      则 $t>1$，且 $x_1=tx_2$。
		      \par\textbf{理由：} $x_1,x_2$ 地位对等，交换二者不影响和与积的结论；设 $t>1$ 可避免绝对值讨论。

		\item[\textbf{步骤二（表达参数 $c$）：}]
		      将 $x_1=tx_2$ 代入已知等式，得
		      \[
			      \begin{aligned}
				      c
				       & =\frac{x_1-x_2}{\ln x_1-\ln x_2}    \\
				       & =\frac{tx_2-x_2}{\ln(tx_2)-\ln x_2} \\
				       & =\frac{x_2(t-1)}{\ln t}.
			      \end{aligned}
		      \]
		      \par\textbf{理由：} 利用 $\ln(tx_2)-\ln x_2=\ln t$，将 $c$ 分离为 $x_2$ 与 $t$ 的表达式。

		\item[\textbf{步骤三（转化目标不等式）：}]
		      先转化和不等式：
		      \[
			      \begin{aligned}
				      x_1+x_2>2c
				       & \iff x_2(t+1)>2\cdot\frac{x_2(t-1)}{\ln t} \\
				       & \iff (t+1)\ln t>2(t-1).
			      \end{aligned}
		      \]
		      再转化积不等式：
		      \[
			      \begin{aligned}
				      x_1x_2<c^2
				       & \iff tx_2^2<\left(\frac{x_2(t-1)}{\ln t}\right)^2 \\
				       & \iff t\ln^2t<(t-1)^2.
			      \end{aligned}
		      \]
		      \par\textbf{理由：} 因为 $x_2>0$ 且 $\ln t>0$，所以可以安全约去正数并保持不等号方向不变。

		\item[\textbf{步骤四（证明和不等式）：}]
		      考虑函数
		      \[
			      F(t)=(t+1)\ln t-2(t-1),\qquad t>1.
		      \]
		      则
		      \[
			      F'(t)=\ln t+\frac{1}{t}-1,
			      \qquad
			      F''(t)=\frac{1}{t}-\frac{1}{t^2}=\frac{t-1}{t^2}.
		      \]
		      当 $t>1$ 时，$F''(t)>0$，所以 $F'(t)$ 在 $(1,+\infty)$ 上单调递增。又
		      \[
			      \lim_{t\to1^+}F'(t)=0,
		      \]
		      故 $F'(t)>0$。因此 $F(t)$ 在 $(1,+\infty)$ 上单调递增，且
		      \[
			      \lim_{t\to1^+}F(t)=0,
		      \]
		      所以 $F(t)>0$，即
		      \[
			      (t+1)\ln t>2(t-1).
		      \]
		      \par\textbf{理由：} $F(t)>0$ 正是和不等式转化后的形式。

		\item[\textbf{步骤五（证明积不等式）：}]
		      考虑函数
		      \[
			      G(t)=(t-1)^2-t\ln^2t,
			      \qquad t>1.
		      \]
		      则
		      \[
			      G'(t)=2(t-1)-\ln^2t-2\ln t.
		      \]
		      继续求二阶导数：
		      \[
			      \begin{aligned}
				      G''(t)
				       & =2-\frac{2\ln t}{t}-\frac{2}{t}     \\
				       & =2\left(1-\frac{\ln t+1}{t}\right).
			      \end{aligned}
		      \]
		      令
		      \[
			      U(t)=1-\frac{\ln t+1}{t}.
		      \]
		      则
		      \[
			      U'(t)=\frac{\ln t}{t^2}>0\qquad(t>1),
		      \]
		      且 $\lim\limits_{t\to1^+}U(t)=0$，所以 $U(t)>0$。于是 $G''(t)>0$，从而 $G'(t)$ 在 $(1,+\infty)$ 上单调递增。又
		      \[
			      \lim_{t\to1^+}G'(t)=0,
		      \]
		      所以 $G'(t)>0$。因此 $G(t)$ 在 $(1,+\infty)$ 上单调递增，且
		      \[
			      \lim_{t\to1^+}G(t)=0,
		      \]
		      所以 $G(t)>0$，即
		      \[
			      t\ln^2t<(t-1)^2.
		      \]
		      \par\textbf{理由：} $G(t)>0$ 正是积不等式转化后的形式。
	\end{description}

	\medskip
	\textbf{\textcolor{CProof}{结论回扣}}

	综合步骤四和步骤五，我们证明了当 $t>1$ 时
	\[
		(t+1)\ln t>2(t-1),
		\qquad
		t\ln^2t<(t-1)^2.
	\]
	因此由步骤三可得
	\[
		x_1+x_2>2c,
		\qquad
		x_1x_2<c^2.
	\]
	又由于 $x_1,x_2$ 的地位对称，原结论对任意满足条件的 $x_1,x_2$ 都成立。
\end{proofbox}